\section{Introduction}
\label{sec:introduction}

Neural networks are algorithms that learn to make predictions from examples. They are used to
identify objects in pictures, to analyze speech, and to aid in controlling vehicles. Neural networks
can, however, be fooled. Szegedy et al.~\cite{MapSource01} showed in 2014 that a change in an image
imperceptible to a human eye can make a network give a completely different output. One famous
example of such perturbation was given by Goodfellow et al.~\cite{MapSource02}: after a tiny change,
a neural network that classified a picture of a panda as a panda with 57.7 percent confidence
classified the same image as a gibbon with 99.3 percent confidence (Figure 1, p.~3). Such changes in
input are called adversarial examples. In applications where mistakes cost a lot, it is important to
know the vulnerability of the networks to such perturbations and to evaluate it properly.

One possible way to formalize the problem is to use a distance measure. Perturb a single input
incrementally. Robustness radius is the largest perturbation magnitude that does not change the
network's decision. A large radius means prediction stability; a small radius means that a small
change will change the prediction. There are two standard ways of estimating this radius. Attack
finds a nearby input that changes the decision; if such an input exists, then the true radius cannot
be larger than the observed distance and provides an upper bound $U$. On the other hand, certificate
proves that no perturbation within a certain distance will change the decision and provides a lower
bound $L$. The true radius $\rstar$ therefore lies between these two bounds.

There are known limitations of both methods. Attack may fail to find an existing change. Carlini and
Wagner~\cite{MapSource03} showed that a defense that had stopped all previous attacks offers no real
protection, and Athalye et al.~\cite{MapSource06} noticed the same problem in seven out of nine
defenses they studied. It follows that a failure of an attack is not a proof of safety. Certificate,
in its turn, suffers from the problem of being possibly too conservative: the bound obtained by
certificate may be much smaller than the distance of the perturbation found by attack. Salman et
al.~\cite{MapSource10} showed that even the best version of a common certificate reports radii 1.5
to 5 times smaller than the distances found by attacks on several MNIST networks. The common problem
of these comparisons is that they compare attack results with certificate bounds; in the case of
disagreement, it is impossible to tell which method is wrong. Exact methods have a potential to
solve this problem. For small networks, a mixed-integer program can find the true radius
$\rstar$~\cite{MapSource12}. Several papers used exact values to validate one side of the bound.
Carlini et al.~\cite{Carlini2017Provable} and Berger et al.~\cite{Berger2026Upper} compared the
proximity of attacks to the true radius but did not compare certificates. Weng et
al.~\cite{Weng2018FastLin} reported a certificate, an attack, and an exact value at once, but only
as averages over two small networks and without splitting the two sources of error.

This paper measures both errors on the same input. The difference between the bounds decomposes into
$U - L = (\rstar - L) + (U - \rstar)$. The first term represents underestimation obtained by the
certificate, and the second term represents overestimation made by the attack. Measuring both parts
of the decomposition for each input allows attributing the gap to the corresponding method. This
paper compares the decomposition on networks with one, two, and three hidden layers with a similar
number of parameters (182, 189, and 186, respectively) using five trained networks of each type and
30 test inputs. This research aims to answer the following question: how do certificate conservatism
and attack suboptimality vary across approximately parameter-matched ReLU classifiers with one, two,
and three hidden layers on a fixed two-dimensional classification task under an L-infinity threat
model? Networks and data are purposely small because exact radii can be computed only on this scale;
therefore, the results cannot be generalized to large or deployed networks.

Section~\ref{sec:literature-review} reviews previous works on attacks, certificates, and exact
methods. Section~\ref{sec:method} describes the data, the networks, and the three estimates.
Section~\ref{sec:results} shows the results, and Section~\ref{sec:discussion} discusses their
implications for failed attacks and the limitations of this study.
